% GENERATED FILE. Edit canonical lessons, then run npm run generate:books.
% canonical-source-hash: fnv1a64:d913bf953f90edf7
% canonical-lessons: JA-C126-shiawase, JA-C126-nemutai, JA-C126-anshin, JA-C126-futsuu, JA-C126-umai

\chapter{Happy, Drowsy, Relieved, Ordinary, Tasty}
\label{ch:ja-happy-drowsy-relieved-ordinary-tasty}
\hlchaptermodality{126}

\begin{chapteropening}
\textbf{By the end of this chapter:} \emph{I can say happy, drowsy, relieved, ordinary and tasty.}

Complete the last lesson of chapter 126: I can say happy, drowsy, relieved, ordinary and tasty.
\end{chapteropening}

\section[\texorpdfstring{\ja{しあわせ} (shiawase)}{しあわせ (shiawase)}]{\ja{しあわせ} (shiawase) — happy}

\label{lesson:JA-C126-shiawase}



\begin{quote}
Before the new one: say the Japanese for very young, then the Japanese for plentiful.
\end{quote}

\subsection*{You'll want to know: \ja{しあわせ}}
\textbf{\ja{しあわせ}} — \emph{shiawase} — \textquotedblleft{}happy\textquotedblright{}.

\begin{grammarlens}[title={What you've built}]
One more word for this chapter.
\end{grammarlens}

\subsection*{Your turn}
\begin{itemize}
\raggedright
  \item \emph{Say it:} \emph{shiawase}
  \item \emph{Say it:} \emph{shiawase}, once more
  \item \emph{Say it:} say \emph{shiawase}
  \item \emph{Recall:} say the Japanese for very young, then the Japanese for plentiful, then say \emph{shiawase} again
\end{itemize}

\begin{tcolorbox}[breakable,colback=teal!4,colframe=teal!35!black,title={Before you move on}]
What does \ja{しあわせ} mean? (\textquotedblleft{}Happy\textquotedblright{}.) Say it once more.
\end{tcolorbox}
\section[\texorpdfstring{\ja{ねむたい} (nemutai)}{ねむたい (nemutai)}]{\ja{ねむたい} (nemutai) — drowsy}

\label{lesson:JA-C126-nemutai}



\begin{quote}
Before the new one: say the Japanese for plentiful, then the Japanese for happy.
\end{quote}

\subsection*{You'll want to know: \ja{ねむたい}}
\textbf{\ja{ねむたい}} — \emph{nemutai} — \textquotedblleft{}drowsy\textquotedblright{}.

\begin{grammarlens}[title={What you've built}]
One more word for this chapter.
\end{grammarlens}

\subsection*{Your turn}
\begin{itemize}
\raggedright
  \item \emph{Say it:} \emph{nemutai}
  \item \emph{Say it:} \emph{nemutai}, once more
  \item \emph{Say it:} say \emph{nemutai}
  \item \emph{Recall:} say the Japanese for plentiful, then the Japanese for happy, then say \emph{nemutai} again
\end{itemize}

\begin{tcolorbox}[breakable,colback=teal!4,colframe=teal!35!black,title={Before you move on}]
What does \ja{ねむたい} mean? (\textquotedblleft{}Drowsy\textquotedblright{}.) Say it once more.
\end{tcolorbox}
\section[\texorpdfstring{\ja{あんしん} (anshin)}{あんしん (anshin)}]{\ja{あんしん} (anshin) — relieved}

\label{lesson:JA-C126-anshin}



\begin{quote}
Before the new one: say the Japanese for happy, then the Japanese for drowsy.
\end{quote}

\subsection*{You'll want to know: \ja{あんしん}}
\textbf{\ja{あんしん}} — \emph{anshin} — \textquotedblleft{}relieved\textquotedblright{}.

\begin{grammarlens}[title={What you've built}]
One more word for this chapter.
\end{grammarlens}

\subsection*{Your turn}
\begin{itemize}
\raggedright
  \item \emph{Say it:} \emph{anshin}
  \item \emph{Say it:} \emph{anshin}, once more
  \item \emph{Say it:} say \emph{anshin}
  \item \emph{Recall:} say the Japanese for happy, then the Japanese for drowsy, then say \emph{anshin} again
\end{itemize}

\begin{tcolorbox}[breakable,colback=teal!4,colframe=teal!35!black,title={Before you move on}]
What does \ja{あんしん} mean? (\textquotedblleft{}Relieved\textquotedblright{}.) Say it once more.
\end{tcolorbox}
\section[\texorpdfstring{\ja{ふつう} (futsuu)}{ふつう (futsuu)}]{\ja{ふつう} (futsuu) — ordinary}

\label{lesson:JA-C126-futsuu}



\begin{quote}
Before the new one: say the Japanese for drowsy, then the Japanese for relieved.
\end{quote}

\subsection*{You'll want to know: \ja{ふつう}}
\textbf{\ja{ふつう}} — \emph{futsuu} — \textquotedblleft{}ordinary\textquotedblright{}.

\begin{grammarlens}[title={What you've built}]
One more word for this chapter.
\end{grammarlens}

\subsection*{Your turn}
\begin{itemize}
\raggedright
  \item \emph{Say it:} \emph{futsuu}
  \item \emph{Say it:} \emph{futsuu}, once more
  \item \emph{Say it:} say \emph{futsuu}
  \item \emph{Recall:} say the Japanese for drowsy, then the Japanese for relieved, then say \emph{futsuu} again
\end{itemize}

\begin{tcolorbox}[breakable,colback=teal!4,colframe=teal!35!black,title={Before you move on}]
What does \ja{ふつう} mean? (\textquotedblleft{}Ordinary\textquotedblright{}.) Say it once more.
\end{tcolorbox}
\section[\texorpdfstring{\ja{うまい} (umai)}{うまい (umai)}]{\ja{うまい} (umai) — tasty; skilful}

\label{lesson:JA-C126-umai}



\begin{quote}
Before the new one: say the Japanese for relieved, then the Japanese for ordinary.
\end{quote}

\subsection*{You'll want to know: \ja{うまい}}
\textbf{\ja{うまい}} — \emph{umai} — \textquotedblleft{}tasty; skilful\textquotedblright{}.

\begin{grammarlens}[title={What you've built}]
One more word for this chapter.
\end{grammarlens}

\subsection*{Your turn}
\begin{itemize}
\raggedright
  \item \emph{Say it:} \emph{umai}
  \item \emph{Say it:} \emph{umai}, once more
  \item \emph{Say it:} say \emph{umai}
  \item \emph{Recall:} say the Japanese for relieved, then the Japanese for ordinary, then say \emph{umai} again
\end{itemize}

\begin{tcolorbox}[breakable,colback=teal!4,colframe=teal!35!black,title={Before you move on}]
What does \ja{うまい} mean? (\textquotedblleft{}Tasty; skilful\textquotedblright{}.) Say it once more.
\end{tcolorbox}
